\documentclass{article}
\setlength{\parindent}{0pt}
\setlength{\parskip}{1em}
\usepackage[margin=1in]{geometry}
\usepackage{natbib}

\begin{document}

\thispagestyle{empty}
Dear Dr. Maderspacher, \hfill 27 July 2026   \newline \vspace{-.75em}

We address a central puzzle in plant biology and phenology that has become increasingly important with higher rates of climate change. Scientists have long recognized that warmer temperatures result in earlier springtime events, such as flowering and insect emergence \citep{Root:2003kl}. But as climates have continued to warm in recent years, two additional patterns have emerged: the rate at which these events have advanced earlier in the calendar year has slowed, and their timing has often become more variable  \citep{fu2015declining,vitasse2022great,chuine2025living}. While researchers have invoked complex models to explain these patterns, we show here that they are predicted from a simple temperature model. Our work solves at once two major puzzles and advances ecological forecasting for plants and animals. 

An extensive and rapidly growing body of work has been revising the long-standing temperature model for biological events to explain patterns of reduced sensitivity and higher variation with warming. This basic model---the thermal sum model (or growing degree day model)---predicts events are triggered once cumulative temperatures reach a threshold. Building on this basic model, research has now yielded hundreds of papers in recent years debating whether photoperiod, winter temperatures, or cold hardiness underlie these patterns \citep[e.g.,][]{fu2015declining,zohner2016ncc,ospreebbms,kovaleskipreprint}. Yet no work has examined whether the patterns are consistent with the thermal sum model, and if so, the implications for climate change.

Here, by working across biology and statistics, we show that these patterns are consistent with the long-standing thermal-sum model. Our model predicts both slower advances and increased variance from additional warming, and we validate these model predictions with experimental data and observational lilac flowering data over the last 70 years across the US. In contrast to predictions using current models of biological events, our framework implies that as springtime events advance with anthropogenic warming, they enter a new temperature regime in which the typical time between events in the same ecosystem can increase from a few days to a month or more.

These results challenge research that has assumed factors other than warm temperatures explain these trends. Further, they imply that commonly used statistical tools that assume linear responses or constant variance may be misspecified and will lead to biased or inaccurate forecasts. Thus, in addition to challenging current methodological approaches, our findings will yield major advances in forecasting and fundamental plant biology: we provide a simple and transparent way to encode the thermal-sum model, on top of which complexity may be added and additional factors better studied.

Our work is interdisciplinary. It builds on biological theory from phenology and on mathematical results that arise in the study of stopped random walks. Moreover, our goal is to reach a broad audience at the intersection of ecology, agriculture, plant biology, climate change, public policy, and statistical methodology. We have therefore simplified the mathematical presentation while retaining the key statistical insight: the same temperature-threshold mechanism that has long been used to explain the timing of spring events also predicts when those events will become less sensitive to warming and less synchronized.

Thank you for your consideration,

Jonathan Auerbach, E.M. Wolkovich, and Andrew Gelman

\newpage 
\bibliographystyle{apalike}
\bibliography{..//desynchronization}

\end{document}

Due to the interdisciplinary nature of our work, we suggest reviewers from both ecology and statistics. From ecology, we suggest the consideration of Susan Harrison (NAS, University of California Davis, spharrison@ucdavis.edu), Nils Stenseth (PNAS, University of Oslo, n.c.stenseth@mn.uio.no), Alan Hastings (PNAS, University of California Davis, amhastings@ucdavis.edu), Donald Ort (PNAS, University of Illinois Urbana-Champaign, d-ort@illinois.edu). Ecology reviewers that work at the intersection of phenology and statistics include Eric Post (University of California Davis, post@ucdavis.edu), Christy Rollinson (The Morton Arboretum, crollinson@mortonarb.org), Josh Gray (North Carolina State University, josh\_gray@ncsu.edu), and James Grace (Concordia University, james.grace@concordia.ca). 

For statistics, we suggest those familiar with our methodological approach such as David Siegmund (NAS, Stanford University, dos@stat.stanford.edu), Michael Baron (American University, baron@american.edu), and Yaakov Malinovsky (University of Maryland, yaakovm@umbc.edu). The latter two recently hosted the 9th International Workshop in Sequential Methodologies at which this work was presented. 

We declare no funding sources or competing interests and would select the CC BY-NC-ND license if accepted.  \newline \vspace{-.75em}